\documentclass[10pt,a4paper]{amsart}
\usepackage{amsmath,amssymb,amsthm,graphicx}
\usepackage{verbatim}
\newenvironment{reference}{\comment}{\endcomment}
%\newenvironment{reference}{}{}
\newenvironment{slogan}{\comment}{\endcomment}
\newenvironment{history}{\comment}{\endcomment}

% For commutative diagrams we use Xy-pic
\usepackage[all]{xy}

% We use 2cell for 2-commutative diagrams.
\xyoption{2cell}
% Theorem environments.
%
\theoremstyle{plain}
\newtheorem{theorem}[subsection]{Theorem}
\newtheorem{proposition}[subsection]{Proposition}
\newtheorem{lemma}[subsection]{Lemma}

\theoremstyle{definition}
\newtheorem{definition}[subsection]{Definition}
\newtheorem{example}[subsection]{Example}
\newtheorem{exercise}[subsection]{Exercise}
\newtheorem{situation}[subsection]{Situation}

\theoremstyle{remark}
\newtheorem{remark}[subsection]{Remark}
\newtheorem{remarks}[subsection]{Remarks}

\numberwithin{equation}{subsection}

% Macros
%
\def\lim{\mathop{\mathrm{lim}}\nolimits}
\def\colim{\mathop{\mathrm{colim}}\nolimits}
\def\Spec{\mathop{\mathrm{Spec}}}
\def\Hom{\mathop{\mathrm{Hom}}\nolimits}
\def\Ext{\mathop{\mathrm{Ext}}\nolimits}
\def\SheafHom{\mathop{\mathcal{H}\!\mathit{om}}\nolimits}
\def\SheafExt{\mathop{\mathcal{E}\!\mathit{xt}}\nolimits}
\def\Sch{\mathit{Sch}}
\def\Mor{\mathop{\mathrm{Mor}}\nolimits}
\def\Ob{\mathop{\mathrm{Ob}}\nolimits}
\def\Sh{\mathop{\mathit{Sh}}\nolimits}
\def\NL{\mathop{N\!L}\nolimits}
\def\CH{\mathop{\mathrm{CH}}\nolimits}
\def\proetale{{pro\text{-}\acute{e}tale}}
\def\etale{{\acute{e}tale}}
\def\QCoh{\mathit{QCoh}}
\def\Ker{\mathop{\mathrm{Ker}}}
\def\Im{\mathop{\mathrm{Im}}}
\def\Coker{\mathop{\mathrm{Coker}}}
\def\Coim{\mathop{\mathrm{Coim}}}

% Boxtimes
%
\DeclareMathSymbol{\boxtimes}{\mathbin}{AMSa}{"02}

%
% Macros for moduli stacks/spaces
%
\def\QCohstack{\mathcal{QC}\!\mathit{oh}}
\def\Cohstack{\mathcal{C}\!\mathit{oh}}
\def\Spacesstack{\mathcal{S}\!\mathit{paces}}
\def\Quotfunctor{\mathrm{Quot}}
\def\Hilbfunctor{\mathrm{Hilb}}
\def\Curvesstack{\mathcal{C}\!\mathit{urves}}
\def\Polarizedstack{\mathcal{P}\!\mathit{olarized}}
\def\Complexesstack{\mathcal{C}\!\mathit{omplexes}}
% \Pic is the operator that assigns to X its picard group, usage \Pic(X)
% \Picardstack_{X/B} denotes the Picard stack of X over B
% \Picardfunctor_{X/B} denotes the Picard functor of X over B
\def\Pic{\mathop{\mathrm{Pic}}\nolimits}
\def\Picardstack{\mathcal{P}\!\mathit{ic}}
\def\Picardfunctor{\mathrm{Pic}}
\def\Deformationcategory{\mathcal{D}\!\mathit{ef}}
\usepackage{cleveref}
\crefname{originallemma}{original lemma}{original lemmas}
\Crefname{originallemma}{Original lemma}{Original lemmas}
\makeatletter
\newlabel{lemma-functor-coherent}{{0FZL}{source}{}{sourcepointer-1}{}}
\newlabel{lemma-functor-coherent@cref}{{[originallemma][0][]0FZL}{[1][0][]source}{}{}{}}
\makeatother
\crefname{originalproposition}{original proposition}{original propositions}
\Crefname{originalproposition}{Original proposition}{Original propositions}
\makeatletter
\newlabel{proposition-gabriel-rosenberg}{{0GPJ}{source}{}{sourcepointer-2}{}}
\newlabel{proposition-gabriel-rosenberg@cref}{{[originalproposition][0][]0GPJ}{[1][0][]source}{}{}{}}
\makeatother
\title{Equivalences of coherent module categories on Noetherian schemes come from isomorphisms and invertible modules}
\author{The Stacks Project Authors}
\date{stacks-project:functors.tex:local-snapshot-20261012; GNU FDL 1.2 or later}
\begin{document}
\maketitle
\noindent\textit{Editorial note.} All delivered mathematical text is the author's own, copied verbatim from the pinned Stacks Project chapter file functors.tex; no statement or proof step was rewritten, reordered or removed.\par\medskip
\noindent\textit{Editorial note.} Span recomputation by label pairing (LANE\_BRIEF 10.2): in this chapter the author's proof environment follows the statement environment as a standalone block; the spec covers the complete environment pair (statement span, proof span) as recorded in the pool candidate.\par\medskip
\noindent\textit{Editorial note.} Occupation scan: combined claim\_key + statement\_anchor + generic-line scan over all registration files (active and withdrawn) found 0 hits for functors.tex (fu-occupied.json, 2026-10-14); consistent with the duty manager's six-check dispatch.\par\medskip
\ifdefined\hypertarget\hypertarget{sourcepointer-1}{}\fi
\noindent\textit{Source reference.} Original lemma 0FZL: sources/stacks-chow/tags.txt, source lines 19956--19956. Cross-chapter standard prerequisite cited by the author inside the counted span; it is used as a pointer and its own proof is not part of this item.\par\medskip
\ifdefined\hypertarget\hypertarget{sourcepointer-2}{}\fi
\noindent\textit{Source reference.} Original proposition 0GPJ: sources/stacks-chow/tags.txt, source lines 20829--20829. Cross-chapter standard prerequisite cited by the author inside the counted span; it is used as a pointer and its own proof is not part of this item.\par\medskip
\begin{lemma}
\label{lemma-equivalence-coherent}
Let $X$ and $Y$ be Noetherian schemes. Let
$F : \textit{Coh}(\mathcal{O}_X) \to \textit{Coh}(\mathcal{O}_Y)$
be an equivalence of categories. Then there is an isomorphism $f : Y \to X$
and an invertible $\mathcal{O}_Y$-module $\mathcal{L}$
such that $F(\mathcal{F}) = f^*\mathcal{F} \otimes \mathcal{L}$.
\end{lemma}

\begin{proof}
By Lemma \ref{lemma-functor-coherent} we obtain a unique functor
$F' : \QCoh(\mathcal{O}_X) \to \QCoh(\mathcal{O}_Y)$ extending $F$.
The same is true for the quasi-inverse of $F$ and by the uniqueness
we conclude that $F'$ is an equivalence. By
Proposition \ref{proposition-gabriel-rosenberg}
we find an isomorphism $f : Y \to X$
and an invertible $\mathcal{O}_Y$-module $\mathcal{L}$
such that $F'(\mathcal{F}) = f^*\mathcal{F} \otimes \mathcal{L}$.
Then $f$ and $\mathcal{L}$ work for $F$ as well.
\end{proof}

\end{document}
